\subsection{An explicit Lang--Weil test for pure primary theta torsion}
\label{sec:lang-weil-theta}

\paragraph{Source scope.}
The Lang quotient construction used here is the one in Poonen,
\emph{Multiples of subvarieties in algebraic groups over finite fields},
IMRN 2005, no.~24, 1487--1498, Lemma~5.1 (author PDF, p.~6).
The recovered \texttt{ANABELIAN\_SUCCESSOR}, A.10--A.11, also uses
finite-field character sums on Abel--Jacobi images.  The following
statement gives explicit constants for the genus-two theta translate
and a sufficient test for a point whose \emph{entire order} is an odd
prime power.  It is unconditional.

\paragraph{Finding A12 (a curve-dependent sufficient test).}
Let $C/\mathbf F_Q$ be a smooth projective geometrically connected
genus-two curve in characteristic two, and assume that $C$ has an
$\mathbf F_Q$-rational point.  Write $X=C^{(1)}$ for its one-fold
relative Frobenius twist,
$J=\operatorname{Pic}^0(X)$, and
\[
 0\longrightarrow\mathcal O_X\longrightarrow F_*\mathcal O_C
 \longrightarrow B\longrightarrow0
\]
for relative Frobenius and its rank-one quotient.  Let
$\Theta=\{L\in J:h^0(X,B\otimes L)>0\}$ denote the Raynaud theta
support.  For any subgroup $H\subseteq G:=J(\mathbf F_Q)$, put
$a=[G:H]$ and
$N_H=\#(\Theta(\mathbf F_Q)\cap H)$.  Then
\[
 N_H\ \ge\
 \max\left\{0,\
 \left\lceil
 \frac{Q+1-\left\lfloor2(a+1)\sqrt Q\right\rfloor}{a}
 \right\rceil\right\}.
 \tag{LW.1}
\]
In particular,
\[
       Q+1-2(a+1)\sqrt Q>a
 \tag{LW.2}
\]
ensures that $\Theta(\mathbf F_Q)\cap H$ contains a nonzero point.
The simpler condition
\[
                       Q>(2a+3)^2
 \tag{LW.3}
\]
implies (LW.2).  If $C$ is ordinary, the origin is not in $\Theta$,
and positivity of the numerator in (LW.1) already suffices.

For an odd prime $\ell$ and $r\ge1$, take
$H=J(\mathbf F_Q)[\ell^r]$.  If (LW.2) holds and
\[
             \langle 2,-1\rangle=(\mathbf Z/\ell^r)^\times,
 \tag{LW.4}
\]
then the prime-power bridge A11 applies: after a finite extension of
constants there is a connected cyclic cover of degree $\ell^s$,
for some $1\le s\le r$, whose final degree-$\ell$ Prym has $2$-rank zero.
Consequently the point-separation conclusion of A11 holds for $C$.
This last conclusion uses A11, whereas (LW.1)--(LW.3) are purely
finite-field counting statements.

\begin{proof}
\emph{1. The theta translate and the Lang quotient.}
Choose $P_0\in X(\mathbf F_Q)$.  The line bundle $B$ has degree one.
A degree-one line bundle on a genus-two curve has a nonzero section
exactly when it is $\mathcal O_X(R)$ for a unique geometric point $R$.
Thus the theta support is the image of the closed embedding
\[
 i_B:X\longrightarrow J,\qquad
 R\longmapsto\mathcal O_X(R)\otimes B^{-1}.
\]
It is the usual Abel--Jacobi embedding
$i_{P_0}(R)=\mathcal O_X(R-P_0)$ translated by the
$\mathbf F_Q$-rational point
$\tau=\mathcal O_X(P_0)\otimes B^{-1}$.
No rational effective divisor representing $B$ is required.

Let $\Phi$ denote the $Q$-power Frobenius endomorphism of $J$.
The homomorphism $\Phi-1$ has differential $-1$, so is an
\'etale isogeny.  Its kernel is the constant group scheme
$G=J(\mathbf F_Q)$.  Factoring through the constant subgroup $H$
gives \'etale isogenies
\[
  J\xrightarrow{\pi_H} J/H\xrightarrow{\psi}J,
  \qquad \psi\pi_H=\Phi-1,
  \qquad \ker\psi=G/H.
 \tag{LW.5}
\]
In particular $\deg\psi=a$.  Form the pullback
\[
                  Y=X\times_{i_B,J,\psi}(J/H).
 \tag{LW.6}
\]
It is an \'etale cover of $X$ of degree $a$.

\emph{2. Geometric connectedness, including the characteristic part.}
We give the argument to make clear that $a$ need not be odd.
Over an algebraic closure, if (LW.6) were disconnected, its
monodromy would be a proper subgroup $A_0$ of the constant abelian
deck group $A=G/H$.  Quotienting by $A_0$ would give a nontrivial
intermediate \'etale isogeny
$\psi':J'\to J$ such that $i_B$ lifts to a morphism $v:X\to J'$.
Indeed the resulting $A/A_0$-torsor on $X$ has trivial geometric
monodromy and hence has a section.
Replace $v$ by $v-v(P_0)$.  It now sends $P_0$ to zero and lifts
$i_{P_0}$.  The universal property of the Jacobian supplies a
homomorphism $u:J\to J'$ with
$u i_{P_0}=v-v(P_0)$.  It follows, again by that universal property,
that $\psi'u=\operatorname{id}_J$.  The two abelian varieties have
the same dimension, so $u$ is an isogeny and
$1=\deg(\psi'u)=\deg\psi'\deg u$, a contradiction.

This is also the full abelian \'etale fundamental-group property
of the Abel--Jacobi map; it applies to characteristic-primary
\'etale covers as well as prime-to-characteristic covers.
See Milne, \emph{Jacobian Varieties}, author version of
12 June 2021, Section~9, Proposition~9.1, p.~30;
its characteristic-primary argument is completed on pp.~31--32 \cite{MilneJV}.
Thus $Y$ is smooth, projective and geometrically connected.
Unramified Riemann--Hurwitz gives
\[
                         g(Y)=a(g(X)-1)+1=a+1.
 \tag{LW.7}
\]

\emph{3. Rational fibres and the exact count.}
For $R\in X(\mathbf F_Q)$, put $x=i_B(R)\in G$.
Choose $y\in J(\overline{\mathbf F}_Q)$ with
$(\Phi-1)y=x$, and put $z=\pi_H(y)\in\psi^{-1}(x)$.
Since $\pi_H$ is defined over $\mathbf F_Q$,
\[
                 \Phi(z)-z=\pi_H(x).
 \tag{LW.8}
\]
Thus $z$ is rational if and only if $x\in H$.
The deck group $G/H$ is constant, so in this case all $a$ points
in the fibre are rational; otherwise none is rational.
Consequently
\[
                           \#Y(\mathbf F_Q)=aN_H.
 \tag{LW.9}
\]
The Weil bound for the curve $Y$ and (LW.7) give
$\#Y(\mathbf F_Q)\ge Q+1-2(a+1)\sqrt Q$.
For the precise cohomological input, see Deligne,
\emph{La conjecture de Weil. I}, formula~(1.5.1), p.~275, and
Theorem~1.6, p.~276 \cite{WeilI}.  Since $\#Y(\mathbf F_Q)$ is an integer,
(LW.9) yields (LW.1).

There is at most one point of $X$ mapping to the origin of $J$.
Condition (LW.2) gives $N_H>1$ and hence a nonzero point of $H$ on
$\Theta$.  If $\sqrt Q>2a+3$, then
$Q-2(a+1)\sqrt Q>\sqrt Q>a-1$, proving (LW.3).
If $C$ is ordinary, Frobenius on $H^1(\mathcal O_C)$ is invertible.
The cohomology sequence defining $B$ therefore gives
$H^0(X,B)=0$, so the origin is not on $\Theta$.

\emph{4. The prime-power consequence.}
A nonzero point in $H=G[\ell^r]$ has exact order $\ell^s$ for some
$1\le s\le r$.  Reduction
$(\mathbf Z/\ell^r)^\times\to(\mathbf Z/\ell^s)^\times$ is surjective.
Thus (LW.4) implies
$\langle2,-1\rangle=(\mathbf Z/\ell^s)^\times$.
Apply A11 to this actual order, allowing the finite constant
extension needed to define the deck action and a chosen point
above the prescribed base point.
\end{proof}

\paragraph{Finding A13 (an ordinary example with order $4159$).}
Let $C/\mathbf F_2$ be the smooth projective model of
\[
                  y^2+y=x+x^{-1}+(x-1)^{-1}.
 \tag{LW.10}
\]
It is ordinary of genus two.  Its Raynaud theta support contains
at least four $\mathbf F_{128}$-rational points of exact order
$4159$.  Moreover
$\langle2,-1\rangle=(\mathbf Z/4159)^\times$.
Thus A11 applies to this ordinary curve with a cyclic cover of
prime degree $4159$.

\begin{proof}
There are exactly three simple poles in (LW.10), at $0,1,\infty$.
The Artin--Schreier genus and $p$-rank formulas give genus two and
$2$-rank two; see Pries--Zhu,
\emph{The $p$-rank stratification of Artin--Schreier curves},
Lemma~2.6, p.~712 \cite{PriesZhu2012}.  The three totally ramified points are rational.
For a polynomial model one may put $h=x(x+1)$ and $Y=hy$, obtaining
\[
               Y^2+(x^2+x)Y=x^5+x^3+x^2+x.
 \tag{LW.11}
\]

There are three points over $\mathbf F_2$, precisely the three
ramified points.  Over
$\mathbf F_4=\mathbf F_2(u)$ with $u^2+u+1=0$, the only nonpole
$x$-values are $u,u^2$.  The right side of (LW.10) at these
values is respectively $u^2,u$, both of trace one to $\mathbf F_2$.
Neither has an Artin--Schreier preimage, so
\[
                      \#C(\mathbf F_2)=\#C(\mathbf F_4)=3.
 \tag{LW.12}
\]
If $\alpha_1,\ldots,\alpha_4$ are the Frobenius eigenvalues, then
$\sum\alpha_i=0$ and $\sum\alpha_i^2=2$.
The functional equation of the genus-two Weil polynomial
therefore gives
\[
                          P(T)=T^4-T^2+4.
 \tag{LW.13}
\]
Write $\beta_1,\beta_2$ for the roots of $T^2-T+4$.
The roots of (LW.13) occur in opposite pairs with squares
$\beta_1,\beta_2$.  For odd $m$ the Jacobian point-count formula
\cite[Theorem~19.1(b)]{MilneAV} gives
\[
 \#J(\mathbf F_{2^m})=(1-\beta_1^m)(1-\beta_2^m)
                   =1-t_m+4^m,
 \quad t_m=\beta_1^m+\beta_2^m.
 \tag{LW.14}
\]
The exact recurrence
$t_0=2,\ t_1=1,\ t_m=t_{m-1}-4t_{m-2}$ yields
\[
                  (t_0,\ldots,t_7)=(2,1,-7,-11,17,61,-7,-251).
\]
Consequently
\[
                    \#J(\mathbf F_{128})=16636=4\cdot4159.
 \tag{LW.15}
\]
Here the relative Frobenius twist has the same model, since the
equation is defined over $\mathbf F_2$.

Trial division by the integers from $2$ through
$\lfloor\sqrt{4159}\rfloor=64$ proves that $4159$ is prime.
The following exact modular calculations certify its order condition:
\[
\begin{gathered}
 2079=3^3\cdot7\cdot11,\qquad 4158=2\cdot2079,\\
 2^{2079}\equiv1,\quad
 2^{693}\equiv2554,\quad
 2^{297}\equiv1245,\quad
 2^{189}\equiv102\pmod{4159}.
\end{gathered}
 \tag{LW.16}
\]
They imply $\operatorname{ord}_{4159}(2)=2079$.
This order is odd, so $-1\notin\langle2\rangle$ and
$\langle2,-1\rangle$ has order $4158$.

In (LW.1), now take $Q=128$ and
$H=J(\mathbf F_{128})[4159]$.  By (LW.15), $|H|=4159$ and $a=4$.
The coarse test already holds because $128>11^2$.
For the integer count,
$113^2<100\cdot128<114^2$ gives
\[
       N_H\ge\left\lceil\frac{129-113}{4}\right\rceil=4.
 \tag{LW.17}
\]
The curve is ordinary, so none of these theta points is the origin;
all have exact order $4159$.  This proves the assertion and
verifies both hypotheses of A11.
\end{proof}

\paragraph{Arithmetic evidence for A13.}
The standard-library certificate
\path{lang_weil_4159_certificate.py} checks the finite-field
counts, recurrence, primality, modular orders and integer bound.
An independent script
\path{verify_ordinary_4159_independent.py} obtains (LW.15) instead
as the determinant of the $11$-by-$11$ Sylvester matrix for
$P(T)$ and $T^7-1$.  These exact arithmetic checks complement
the proof above; they do not constitute formal proof-assistant
verification.

\paragraph{Finding A14 (a fixed-prime barrier to the sufficient test).}
Fix a genus-two curve over $\mathbf F_q$ and an odd prime
$\ell\ne\operatorname{char}\mathbf F_q$.  Let $\Phi$ be Frobenius on
the rank-four module $T_\ell J$, and choose $d\ge1$ with
$\Phi^d\equiv1\pmod\ell$.  One may take its order modulo $\ell$,
or the explicit integer
$d=|\operatorname{GL}_4(\mathbf F_\ell)|
 =\prod_{j=0}^3(\ell^4-\ell^j)$.
Set
\[
                 C_\ell=\ell^{v_\ell(\#J(\mathbf F_{q^d}))}.
 \tag{LW.18}
\]
For every $m\ge1$,
\[
       |J(\mathbf F_{q^m})[\ell^\infty]|
                    \le C_\ell\,\ell^{4v_\ell(m)}
                    \le C_\ell m^4.
 \tag{LW.19}
\]
Let $M$ be the first integer $m\ge4$ such that
\[
                         q^{3m}\ge4C_\ell^2m^8.
 \tag{LW.20}
\]
This is a finite, computable integer.  For every $m\ge M$ and
every $r\ge1$, condition (LW.2) fails for
$Q=q^m$ and $H=J(\mathbf F_{q^m})[\ell^r]$.
The assertion concerns the sufficient numerical test, and does
not assert absence of theta torsion.

\begin{proof}
Put $U=\Phi^d-1\in\ell M_4(\mathbf Z_\ell)$.
For every positive integer $t$, the binomial identity gives
\[
 \Phi^{dt}-1=tU V_t,\qquad
 V_t=1+\sum_{j=2}^t\frac{\binom tj}{t}U^{j-1}.
 \tag{LW.21}
\]
Since
$\binom tj/t=\binom{t-1}{j-1}/j$ and
$j-1-v_\ell(j)\ge1$ for odd $\ell$ and $j\ge2$,
every summand after $1$ lies in $\ell M_4(\mathbf Z_\ell)$.
Thus $V_t\equiv1\pmod\ell$ and its determinant is a unit.
No eigenvalue of $\Phi^d$ is one, by the Weil absolute values.
Taking determinants and valuations in (LW.21) proves the exact formula
\[
 v_\ell\#J(\mathbf F_{q^{dt}})
   =v_\ell\#J(\mathbf F_{q^d})+4v_\ell(t).
 \tag{LW.22}
\]
The inclusion
$J(\mathbf F_{q^m})\subseteq J(\mathbf F_{q^{dm}})$
and (LW.22) with $t=m$ prove (LW.19).

Write $s=q^{m/2}$ and
$a=[J(\mathbf F_{q^m}):J(\mathbf F_{q^m})[\ell^r]]$.
The Weil lower bound for the abelian surface and (LW.19) imply
\[
                     a\ge\frac{(s-1)^4}{C_\ell m^4}.
 \tag{LW.23}
\]
Condition (LW.2), after rearrangement, would require
\[
                        a<\frac{(s-1)^2}{2s+1}.
 \tag{LW.24}
\]
It is therefore impossible whenever
$(s-1)^2(2s+1)\ge C_\ell m^4$.
For $m\ge4$ we have $s\ge2$, and
$(s-1)^2(2s+1)\ge s^3/2$.
Condition (LW.20) gives $s^3/2\ge C_\ell m^4$,
proving failure of (LW.2) at that $m$.
Finally $q^{3m}/m^8$ tends to infinity and is strictly increasing
for integral $m\ge4$: its consecutive ratio is at least
$8(4/5)^8>1$.  Hence $M$ exists and the same conclusion holds
for every $m\ge M$.
\end{proof}

\paragraph{What an analytic prime theorem still does not supply.}
A12 allows a curve-dependent successful test and A13 proves that
it can succeed on an ordinary curve.  A14 rules out making this
particular test succeed automatically by enlarging constants at one
fixed odd prime.  Allowing the prime to vary would require additional
information: a sufficiently large pure prime-power subgroup of
$J(\mathbf F_{q^m})$ together with the unit-group condition (LW.4).
For example, (LW.2) requires its order to dominate the cofactor on
the scale specified by (LW.24).
The Jacobian orders are the integers
$\det(1-\Phi^m)$; a density theorem for primes at which $2$ has a
specified multiplicative order does not by itself control large
prime-power divisors of this sequence.
This is not asserted to be the strongest numerical criterion for
point separation.  The recovered \texttt{ANABELIAN\_SUCCESSOR},
A.10, gives a different rational criterion: for a quotient
$\chi:J_C(\mathbf F_Q)\twoheadrightarrow\mathbf Z/\ell$ with
kernel $K$, writing $N=\#C(\mathbf F_Q)$, the conditions
$N\ge2|K|$ and $\ell>N$, together with $2$ primitive modulo $\ell$,
give a suitable rational prime-degree cover.  That criterion can
permit larger cofactors and does not require pure-primary theta
incidence.  Its order hypothesis is stronger than (LW.4):
the prime $4159$ of A13 has $2$ of index two, not a primitive root.
The factorization and order conditions must therefore be stated
for the particular criterion being used.
Once the finite Lang quotient is specified, its rational points
are already controlled by the unconditional finite-field Weil bound.
Thus no Dedekind-zeta GRH, Artin holomorphy assumption, or
number-field least-prime theorem is a premise of A12--A14.
An inference that such an analytic theorem forces the missing
geometric incidence or the needed factorization would require a
separate argument.
